% Fuente única de clase: Recursos/2025-08-12.pdf, páginas 1--9.
En la clase 1 fijamos el lenguaje de conjuntos y funciones y lo usamos para introducir categorías. La inyectividad y la sobreyectividad admiten formulaciones mediante inversas y cancelación; la factorización de una función por sus fibras conecta las relaciones de equivalencia con los cocientes.

Los fragmentos originales de \citet{aluffi2009} incorporados a las pizarras contienen la definición I.1.1, las proposiciones I.2.1 e I.2.3, el teorema I.2.7 y la definición I.3.1. Sirven de orientación para las equivalencias y construcciones de esta clase; las demostraciones siguientes desarrollan también los pasos que las pizarras dejan indicados.

%%%%%%%%%%%%%%%%%%%%%%%%%%%%%%%%%%%%%%%%%%%%%%%%%%%%%%%%%%%%%%%%%%%%%%%%%%%%%%%
\section{Conjuntos, lógica y relaciones}
%%%%%%%%%%%%%%%%%%%%%%%%%%%%%%%%%%%%%%%%%%%%%%%%%%%%%%%%%%%%%%%%%%%%%%%%%%%%%%%
%%%%%%%%%%%%%%%%%%%%%%%%%%%%%%%%%%%%%%%%%%%%%%%%%%%%%%%%%%%%%%%%%%%%%%%%%%%%%%%
\subsection{Descripción de conjuntos y cuantificadores}
%%%%%%%%%%%%%%%%%%%%%%%%%%%%%%%%%%%%%%%%%%%%%%%%%%%%%%%%%%%%%%%%%%%%%%%%%%%%%%%
% P01--P02.
Un conjunto puede describirse por \emph{extensión}, enumerando sus elementos, o por \emph{comprensión}, mediante una condición que los caracteriza. Por ejemplo,
\[
\set{1,2,3,4}=\set{x\in\mathbb Z\given1\le x\le4},
\qquad \mathbb Z=\{\ldots,-2,-1,0,1,2,\ldots\}.
\]
En una expresión $\set{x\in U\given P(x)}$ fijamos un conjunto ambiente $U$ y seleccionamos los elementos que satisfacen el predicado $P$.

Los conectivos lógicos $\land$, $\lor$, $\Rightarrow$, $\neg$ y $\Leftrightarrow$ expresan, respectivamente, conjunción, disyunción, implicación, negación y equivalencia. La disyunción es inclusiva. Los cuantificadores $\forall$ y $\exists$ significan «para todo» y «existe». Las abreviaturas con dominio se interpretan como
\[
\begin{aligned}
(\forall x\in U)\,P(x)&\quad\Longleftrightarrow\quad
\forall x\,(x\in U\Rightarrow P(x)),\\
(\exists x\in U)\,P(x)&\quad\Longleftrightarrow\quad
\exists x\,(x\in U\land P(x)).
\end{aligned}
\]
En particular, negar un enunciado universal equivale a encontrar un contraejemplo:
\[
\neg\bigl((\forall x\in U)P(x)\bigr)
\quad\Longleftrightarrow\quad (\exists x\in U)\neg P(x).
\]

\begin{ejem}
El enunciado «si $x$ es entero, entonces $x$ es racional y positivo» se escribe
\[
(\forall x\in\mathbb Z)\,(x\in\mathbb Q\land x>0).
\]
Es falso: $-2$ es entero y no es positivo. En cambio, todo entero es racional. En efecto, si $x\in\mathbb Z$ es arbitrario, entonces $x=x/1$ es una razón de enteros con denominador no nulo; por tanto, $x\in\mathbb Q$.
\end{ejem}
La segunda demostración ilustra la prueba de una implicación universal: fijamos un elemento arbitrario, suponemos la hipótesis y deducimos la conclusión.

%%%%%%%%%%%%%%%%%%%%%%%%%%%%%%%%%%%%%%%%%%%%%%%%%%%%%%%%%%%%%%%%%%%%%%%%%%%%%%%
\subsection{Operaciones y productos}
%%%%%%%%%%%%%%%%%%%%%%%%%%%%%%%%%%%%%%%%%%%%%%%%%%%%%%%%%%%%%%%%%%%%%%%%%%%%%%%
% P03.
Si $A,B\subseteq U$, definimos
\[
A\cup B=\set{x\in U\given x\in A\lor x\in B},\qquad
A\cap B=\set{x\in U\given x\in A\land x\in B}.
\]
Para una familia $(A_i)_{i\in I}$ de subconjuntos de $U$,
\[
\begin{aligned}
\bigcup_{i\in I}A_i&=\set{x\in U\given(\exists i\in I)\ x\in A_i},\\
\bigcap_{i\in I}A_i&=\set{x\in U\given(\forall i\in I)\ x\in A_i}.
\end{aligned}
\]
El producto cartesiano registra elecciones ordenadas:
\[
A\times B=\set{(a,b)\given a\in A,\ b\in B}.
\]
Un par $(a,b)$ también puede verse como la función $f:\{1,2\}\to A\cup B$ con $f(1)=a$ y $f(2)=b$. Generalizando esta interpretación,
\[
\prod_{i\in I}A_i=
\set*{f:I\longrightarrow\bigcup_{i\in I}A_i\given
                 (\forall i\in I)\ f(i)\in A_i}.
\]
Un elemento del producto elige un elemento de cada factor y conserva su índice.

\begin{nota}
En la página 3 se alternan las letras $A_i$ y $X_i$ en la fórmula del producto. Usamos $A_i$ en toda la expresión para referirnos a una misma familia. Éste es únicamente un ajuste de notación.
\end{nota}

%%%%%%%%%%%%%%%%%%%%%%%%%%%%%%%%%%%%%%%%%%%%%%%%%%%%%%%%%%%%%%%%%%%%%%%%%%%%%%%
\subsection{Equivalencias y conjuntos cociente}
%%%%%%%%%%%%%%%%%%%%%%%%%%%%%%%%%%%%%%%%%%%%%%%%%%%%%%%%%%%%%%%%%%%%%%%%%%%%%%%
% P03--P04.
Una \emph{relación} sobre un conjunto $S$ es un subconjunto $R\subseteq S\times S$. Escribimos $a\sim_R b$ cuando $(a,b)\in R$.
\begin{defi}
Una relación $\sim$ en $S$ es de \emph{equivalencia} si es reflexiva, simétrica y transitiva; esto es, para todos $a,b,c\in S$,
\[
a\sim a,\qquad a\sim b\Rightarrow b\sim a,
\qquad (a\sim b\land b\sim c)\Rightarrow a\sim c.
\]
La \emph{clase de equivalencia} de $a$ y el \emph{conjunto cociente} son
\[
[a]_R=\set{b\in S\given b\sim_R a},\qquad
S/{\sim_R}=\set{[a]_R\given a\in S}.
\]
\end{defi}
La simetría permite escribir indistintamente $b\sim a$ o $a\sim b$ en la definición de $[a]$. Dos clases son iguales o disjuntas: si $c\in[a]\cap[b]$, entonces $a\sim c\sim b$, y la transitividad muestra que todo elemento relacionado con $a$ también está relacionado con $b$, y recíprocamente.

%%%%%%%%%%%%%%%%%%%%%%%%%%%%%%%%%%%%%%%%%%%%%%%%%%%%%%%%%%%%%%%%%%%%%%%%%%%%%%%
\section{Funciones e inversas}
%%%%%%%%%%%%%%%%%%%%%%%%%%%%%%%%%%%%%%%%%%%%%%%%%%%%%%%%%%%%%%%%%%%%%%%%%%%%%%%
%%%%%%%%%%%%%%%%%%%%%%%%%%%%%%%%%%%%%%%%%%%%%%%%%%%%%%%%%%%%%%%%%%%%%%%%%%%%%%%
\subsection{Funciones, composición e inyectividad}
%%%%%%%%%%%%%%%%%%%%%%%%%%%%%%%%%%%%%%%%%%%%%%%%%%%%%%%%%%%%%%%%%%%%%%%%%%%%%%%
% P04--P05.
\begin{defi}
Una función $f:A\to B$ es una relación $f\subseteq A\times B$ tal que para cada $a\in A$ existe un único $b\in B$ con $(a,b)\in f$. Este elemento se denota $f(a)$. El conjunto de funciones de $A$ en $B$ se denota $B^A$.
\end{defi}
Podemos especificar una función mediante su regla de asignación. Por ejemplo,
\[
|\cdot|:\mathbb R\to\mathbb R,\qquad
|x|=\begin{cases}x,&x\ge0,\\-x,&x<0.\end{cases}
\]
Para $f:A\to B$ y $g:B\to C$, su composición es
\[
g\circ f:A\to C,\qquad (g\circ f)(a)=g(f(a)).
\]
Es asociativa, pues ambas funciones $h\circ(g\circ f)$ y $(h\circ g)\circ f$ envían $a$ a $h(g(f(a)))$. La identidad $1_A$ envía cada $a$ a sí mismo.

\begin{defi}
La función $f:A\to B$ es \emph{inyectiva} si $f(a)=f(a')$ implica $a=a'$; equivalentemente, $a\ne a'$ implica $f(a)\ne f(a')$. Es \emph{sobreyectiva} si para cada $b\in B$ existe $a\in A$ con $f(a)=b$, es decir, $\im f=B$. Es \emph{biyectiva} si satisface ambas condiciones.
\end{defi}

\begin{nota}
En la página 4 se formula primero sólo la unicidad de la imagen y luego se escribe la condición de existencia y unicidad. La unicidad aislada permite relaciones sin imagen para algunos elementos; la definición de función requiere ambas condiciones.
\end{nota}

%%%%%%%%%%%%%%%%%%%%%%%%%%%%%%%%%%%%%%%%%%%%%%%%%%%%%%%%%%%%%%%%%%%%%%%%%%%%%%%
\subsection{Inversas por un lado}
%%%%%%%%%%%%%%%%%%%%%%%%%%%%%%%%%%%%%%%%%%%%%%%%%%%%%%%%%%%%%%%%%%%%%%%%%%%%%%%
\begin{prop}\label{prop:C1-inversas}
Sea $f:A\to B$, con $A\ne\varnothing$.
\begin{enumerate}[label=\alph*)]
\item Existe $g:B\to A$ con $g\circ f=1_A$ si y sólo si $f$ es inyectiva.
\item Existe $s:B\to A$ con $f\circ s=1_B$ si y sólo si $f$ es sobreyectiva, usando el axioma de elección para la implicación de derecha a izquierda.
\end{enumerate}
\end{prop}
\begin{proof}
Si $g\circ f=1_A$ y $f(a)=f(a')$, aplicamos $g$ y obtenemos $a=a'$. Recíprocamente, si $f$ es inyectiva, fijamos $a_0\in A$ y definimos $g(b)$ como el único $a$ con $f(a)=b$ cuando $b\in\im f$, y como $a_0$ cuando $b\notin\im f$. Entonces $g(f(a))=a$.

Si $f\circ s=1_B$, para todo $b\in B$ tenemos $b=f(s(b))$, de modo que $f$ es sobreyectiva. A la inversa, si todas las fibras $f^{-1}(\{b\})$ son no vacías, elegimos en cada una un elemento $s(b)$; esta elección satisface $f(s(b))=b$.
\end{proof}

\begin{obs}
La hipótesis $A\ne\varnothing$ importa en el primer inciso: la función $\varnothing\to\{1\}$ es inyectiva, pero no tiene inversa izquierda porque no existe función $\{1\}\to\varnothing$. Es el ejemplo de la página 5. El segundo inciso sigue siendo válido cuando $A=\varnothing$: una función de ese dominio es sobreyectiva exactamente cuando también $B=\varnothing$.
\end{obs}
\begin{nota}
Las pizarras reproducen el enunciado sobre inversas, pero no desarrollan su prueba. La construcción anterior completa ese paso y hace explícito el uso de elección en la selección de preimágenes. Para una biyección no se necesita elegir entre varias: cada preimagen es única y determina su inversa bilateral.
\end{nota}

%%%%%%%%%%%%%%%%%%%%%%%%%%%%%%%%%%%%%%%%%%%%%%%%%%%%%%%%%%%%%%%%%%%%%%%%%%%%%%%
\subsection{Monomorfismos y epimorfismos de conjuntos}
%%%%%%%%%%%%%%%%%%%%%%%%%%%%%%%%%%%%%%%%%%%%%%%%%%%%%%%%%%%%%%%%%%%%%%%%%%%%%%%
% P06--P07.
\begin{defi}
Una función $f:A\to B$ es un \emph{monomorfismo} si para todo conjunto $Z$ y cualesquiera $\alpha',\alpha'':Z\to A$,
\[
f\circ\alpha'=f\circ\alpha''\quad\Rightarrow\quad\alpha'=\alpha''.
\]
Es un \emph{epimorfismo} si para todo conjunto $Z$ y cualesquiera $\beta',\beta'':B\to Z$,
\[
\beta'\circ f=\beta''\circ f\quad\Rightarrow\quad\beta'=\beta''.
\]
\end{defi}

\begin{prop}\label{prop:C1-monoepi}
En conjuntos, los monomorfismos son exactamente las funciones inyectivas y los epimorfismos son exactamente las funciones sobreyectivas.
\end{prop}
\begin{proof}
Supongamos primero $A\ne\varnothing$ y $f$ inyectiva. Por la proposición~\ref{prop:C1-inversas} tiene inversa izquierda $g$. Si $f\alpha'=f\alpha''$, entonces
\[
\alpha'=g(f\alpha')=g(f\alpha'')=\alpha''.
\]
Si $A=\varnothing$, siempre hay a lo sumo una función $Z\to A$, y la cancelación también se cumple. Recíprocamente, si $f$ es mono y $f(a)=f(a')$, las funciones $\{*\}\to A$ que toman los valores $a$ y $a'$ tienen el mismo compuesto con $f$. Deben ser iguales; luego $a=a'$.

Para la segunda equivalencia usamos la estrategia de la pizarra. Si $f$ no es sobreyectiva, fijamos $b_0\in B\setminus\im f$ y definimos
\[
\beta'(b)=0,\qquad
\beta''(b)=\begin{cases}1,&b=b_0,\\0,&b\ne b_0.\end{cases}
\]
Son funciones $B\to\{0,1\}$ distintas con $\beta'f=\beta''f$. Por tanto, $f$ no es epi.

Si $f$ es sobreyectiva y $A\ne\varnothing$, una inversa derecha $s$ permite cancelar:
\[
\beta'=(\beta'f)s=(\beta''f)s=\beta''.
\]
Si $A=\varnothing$, la sobreyectividad obliga a $B=\varnothing$ y sólo existe una función $B\to Z$. Esto termina la prueba. También puede evitarse elección en esta última dirección: para cada $b$ usamos cualquier $a$ con $f(a)=b$ y obtenemos $\beta'(b)=\beta''(b)$.
\end{proof}

%%%%%%%%%%%%%%%%%%%%%%%%%%%%%%%%%%%%%%%%%%%%%%%%%%%%%%%%%%%%%%%%%%%%%%%%%%%%%%%
\subsection{Factorización por las fibras}
%%%%%%%%%%%%%%%%%%%%%%%%%%%%%%%%%%%%%%%%%%%%%%%%%%%%%%%%%%%%%%%%%%%%%%%%%%%%%%%
% P07: el recorte enuncia «como arriba»; aquí se repone la relación necesaria.
Para una función $f:A\to B$, definimos $a\sim a'$ si $f(a)=f(a')$. La igualdad de imágenes hace que esta relación sea reflexiva, simétrica y transitiva.
\begin{teo}\label{teo:C1-factorizacion}
Toda función se factoriza como una sobreyección, una biyección y una inclusión:
\[
\begin{tikzcd}[column sep=large]
A\arrow[r,two heads,"q"]\arrow[rrr,bend left=25,"f"]
&A/{\sim}\arrow[r,"\widetilde f","\simeq"']
&\im f\arrow[r,hook,"i"]&B.
\end{tikzcd}
\]
Aquí $q(a)=[a]$ y $\widetilde f([a])=f(a)$.
\end{teo}
\begin{proof}
La proyección $q$ es sobreyectiva por definición del cociente. La regla de $\widetilde f$ está bien definida porque representantes equivalentes tienen igual imagen. Es sobreyectiva sobre $\im f$; es inyectiva porque $f(a)=f(a')$ implica $[a]=[a']$. La inclusión $i$ es inyectiva y $i\widetilde f q(a)=f(a)$ para todo $a$.
\end{proof}
\begin{nota}
El teorema reproducido en la página 7 remite a una relación definida «arriba», fuera del fragmento visible. Hemos explicitado la relación de igualdad de imágenes y comprobado la factorización; no se presupone una explicación oral ausente.
\end{nota}

%%%%%%%%%%%%%%%%%%%%%%%%%%%%%%%%%%%%%%%%%%%%%%%%%%%%%%%%%%%%%%%%%%%%%%%%%%%%%%%
\section{Categorías y grupoides}
%%%%%%%%%%%%%%%%%%%%%%%%%%%%%%%%%%%%%%%%%%%%%%%%%%%%%%%%%%%%%%%%%%%%%%%%%%%%%%%
%%%%%%%%%%%%%%%%%%%%%%%%%%%%%%%%%%%%%%%%%%%%%%%%%%%%%%%%%%%%%%%%%%%%%%%%%%%%%%%
\subsection{Objetos, morfismos e isomorfismos}
%%%%%%%%%%%%%%%%%%%%%%%%%%%%%%%%%%%%%%%%%%%%%%%%%%%%%%%%%%%%%%%%%%%%%%%%%%%%%%%
% P07--P08.
\begin{defi}
Una \emph{categoría} $\mathcal C$ consta de una clase $\Ob(\mathcal C)$ de objetos, un conjunto $\Hom_{\mathcal C}(A,B)$ de morfismos para cada par de objetos y una composición
\[
\Hom_{\mathcal C}(A,B)\times\Hom_{\mathcal C}(B,C)
\longrightarrow\Hom_{\mathcal C}(A,C),\qquad(f,g)\longmapsto gf.
\]
La composición es asociativa: $(hg)f=h(gf)$. Cada objeto $A$ tiene un morfismo identidad $1_A$ tal que $f1_A=f$ y $1_Bf=f$ para todo $f:A\to B$.
\end{defi}
El ejemplo básico es $\Set$, con conjuntos como objetos y funciones como morfismos.

\begin{defi}
Un morfismo $f:A\to B$ es un \emph{isomorfismo} si existe $g:B\to A$ con $gf=1_A$ y $fg=1_B$. Escribimos $A\cong B$ cuando existe tal isomorfismo. Un \emph{grupoide} es una categoría cuyos morfismos son todos isomorfismos.
\end{defi}
\[
\begin{tikzcd}[column sep=large]
A\arrow[r,bend left=20,"f"]&B\arrow[l,bend left=20,"g"]
\end{tikzcd}
\qquad gf=1_A,\quad fg=1_B.
\]
La inversa es única: si $g$ y $h$ son inversas de $f$, entonces $g=g(fh)=(gf)h=h$.

%%%%%%%%%%%%%%%%%%%%%%%%%%%%%%%%%%%%%%%%%%%%%%%%%%%%%%%%%%%%%%%%%%%%%%%%%%%%%%%
\subsection{La categoría de una relación}
%%%%%%%%%%%%%%%%%%%%%%%%%%%%%%%%%%%%%%%%%%%%%%%%%%%%%%%%%%%%%%%%%%%%%%%%%%%%%%%
% P08--P09.
Sea $\sim$ una relación reflexiva y transitiva sobre $S$. Definimos una categoría $\mathcal C$ con objetos los elementos de $S$ y
\[
\Hom_{\mathcal C}(a,b)=
\begin{cases}\{(a,b)\},&a\sim b,\\\varnothing,&a\not\sim b.\end{cases}
\]
La reflexividad proporciona $1_a=(a,a)$. Si $f=(a,b)$ y $g=(b,c)$ existen, la transitividad garantiza que $(a,c)$ también existe, y definimos $gf=(a,c)$. Las identidades actúan como tales; la asociatividad se sigue porque, entre dos objetos fijos, hay a lo sumo un morfismo.

Si además $\sim$ es simétrica, todo $(a,b)$ tiene inversa $(b,a)$, puesto que
\[
(b,a)(a,b)=(a,a),\qquad (a,b)(b,a)=(b,b).
\]
Así, una relación de equivalencia determina un grupoide. Cuando la relación es la igualdad, sólo existen las identidades:
\[
\Hom_{\mathcal C}(a,b)=
\begin{cases}\{1_a\},&a=b,\\\varnothing,&a\ne b.\end{cases}
\]
Ésta es la \emph{categoría discreta} sobre $S$. En la clase 2 usaremos los grupoides para interpretar los axiomas de grupo.
